\section{Introduction}

In vitro fertilization (IVF) is now a widely used treatment for infertility, and the quality of the single embryo selected for transfer is a primary determinant of clinical success and of the risk of unsafe multiple pregnancies. Embryo selection is strongly influenced by ploidy status: euploid embryos, with a normal chromosomal complement, tend to yield successful pregnancies, whereas aneuploid embryos are associated with failed implantation, miscarriage, and chromosomal disorders. Aneuploidy risk rises sharply with advancing maternal age, and the relationship between age and embryonic aneuploidy has been documented in large comprehensive chromosomal screening studies \citep{franasiak2014aneuploidy}.

The current reference standard for chromosomal assessment is preimplantation genetic testing for aneuploidy (PGT-A), which requires a trophectoderm biopsy, whole-genome amplification, and copy-number analysis using next-generation sequencing (NGS) or related platforms \citep{garciapascual2020ngs,friedenthal2020ngs}. PGT-A is invasive, costly, and can be confounded by embryonic mosaicism --- the presence of both euploid and aneuploid cells within a biopsy --- which contributes to discordant diagnoses and, in some cases, to miscarriage of embryos initially classified as euploid \citep{maxwell2016miscarriage}. These limitations motivate complementary, non-invasive methods for ploidy prediction.

The widespread adoption of time-lapse imaging systems in embryology laboratories, together with the growth of associated image and outcome datasets, has enabled the development of deep-learning models for automated embryo assessment. STORK \citep{khosravi2019deep} used a convolutional neural network trained on time-lapse images to predict blastocyst quality, with performance comparable to individual embryologists. Subsequent models explicitly target ploidy: ERICA \citep{chavezbadiola2020erica} analyzed 1{,}231 static blastocyst images to predict ploidy status with 70\% accuracy, an area under the receiver operating characteristic curve (AUC) of 74\%, a sensitivity of 54\%, and a specificity of 86\%, and placed a euploid embryo in the top rank in 79\% of cases. Barnes et al.\ developed machine-learning models that combine a single image at 110 hours post insemination (hpi) with morphokinetic and clinical variables to predict ploidy \citep{barnes2023nonhinvasive}. Lee et al.\ trained an end-to-end deep learning model on 690 full time-lapse video sequences and reported an AUC of 0.74 for distinguishing aneuploid embryos from euploid and mosaic embryos \citep{lee2021endtoend}. More recent studies have continued to explore convolutional, recurrent, and transformer architectures for ploidy classification from time-lapse data \citep{paya2023deeplearning,chen2023knowledge,handayani2024unet,canat2024morphokinetic,avidiansyah2025vit,maekawa2026aiprediction,ma2024idascore}, and comprehensive benchmarks have shown that even the best models remain limited by data volume and class imbalance \citep{bamford2023ml12,zou2022timelapse}.

A key challenge in analyzing time-lapse sequences for ploidy prediction is that not all developmental stages carry equally informative signal. Prior work has therefore focused on feature extraction from specific windows of development, such as the timing and presence of blastocyst expansion on day 5, although the predictive accuracy of such criteria has been shown to vary across clinics. Using the full time-lapse video circumvents manual time-point selection but can introduce noise and requires larger training sets. In this work, we address these challenges by introducing BELA (Blastocyst Evaluation Learning Algorithm), a fully automated ploidy prediction model that requires only day-5 time-lapse frames (96--112 hpi) and maternal age as inputs. BELA first predicts a blastocyst score from time-lapse features using a multitask bidirectional LSTM and then uses this model-derived score, optionally combined with maternal age, to predict ploidy using logistic regression. The contributions of this paper are: (i) a two-step automated pipeline that obviates subjective embryologist annotation while matching or exceeding the performance of embryologist-scored baselines, (ii) multitask learning that jointly predicts inner cell mass (ICM), trophectoderm (TE), expansion, and overall blastocyst score, (iii) extensive internal and external validation on four datasets spanning two instruments and three clinics (WCM-Embryoscope, WCM-Embryoscope+, IVI Valencia, and IVF Florida), and (iv) a web-based clinical tool, STORK-V, that deploys BELA for use in embryology laboratories.

\section{Related Work}

\paragraph{Deep learning for embryo quality assessment.}
Automated assessment of embryo morphology from time-lapse imaging has been explored by multiple groups \citep{khosravi2019deep,paya2022contrastive,canat2024morphokinetic,jiang2023airreview}. STORK \citep{khosravi2019deep} demonstrated that a CNN-based scoring model could predict blastocyst quality consistent with birth outcomes, and subsequent work has extended these ideas to different developmental stages and imaging platforms. Review articles have surveyed the expanding use of artificial intelligence across the IVF laboratory \citep{jiang2023airreview}.

\paragraph{Ploidy prediction from images and video.}
Several studies have specifically targeted ploidy status. ERICA \citep{chavezbadiola2020erica} used static images of blastocysts to predict euploidy with reported accuracy of 70\% and a top-rank euploid-selection rate of 79\%. STORK-A \citep{barnes2023nonhinvasive} combined a single 110~hpi image with morphokinetic and clinical features. Lee et al.\ \citep{lee2021endtoend} and Paya et al.\ \citep{paya2023deeplearning} investigated end-to-end video classification with CNN-BiLSTM architectures, while Chen et al.\ \citep{chen2023knowledge} integrated multi-focus video and clinical information for euploidy prediction in couples with chromosomal rearrangements. Bamford et al.\ \citep{bamford2023ml12} compared twelve machine-learning models on a morphokinetic meta-dataset and found that mixed-effects logistic regression remained competitive. More recent work has explored U-Net segmentation \citep{handayani2024unet}, vision transformers \citep{avidiansyah2025vit}, spatiotemporal CNNs \citep{maekawa2026aiprediction}, and large-scale foundation models for IVF imaging \citep{rajendran2025femi}. Commercial scoring systems such as iDAScore \citep{ma2024idascore} have also been evaluated for non-invasive aneuploidy prediction.

\paragraph{Inter-observer variability of manual grading.}
A parallel motivation for automated scoring is the substantial inter- and intra-observer variability of manual blastocyst grading. Multiple studies have reported only moderate agreement between embryologists on grading decisions \citep{paternot2009intra,paternot2011multicentre,hammond2020freeze}, and simplified grading schemes have been proposed to improve consistency \citep{richardson2015grading}. These observations reinforce the need for objective, reproducible, and standardized embryo assessment tools.

\paragraph{Distinction from prior work.}
Unlike models relying on embryologist-derived scores or full-length video processing, BELA automatically predicts a blastocyst score from a narrowly defined day-5 window (96--112~hpi) using a multitask BiLSTM and then uses this model-derived score as an interpretable intermediate representation for ploidy classification. BELA therefore combines the data efficiency of targeted temporal windows with the interpretability of a score-based intermediate, while dispensing with subjective embryologist annotation.
